% Options for packages loaded elsewhere
\PassOptionsToPackage{unicode}{hyperref}
\PassOptionsToPackage{hyphens}{url}
%
\documentclass[
  10pt,
]{article}
\usepackage{amsmath,amssymb}
\usepackage{iftex}
\ifPDFTeX
  \usepackage[T1]{fontenc}
  \usepackage[utf8]{inputenc}
  \usepackage{textcomp} % provide euro and other symbols
\else % if luatex or xetex
  \usepackage{unicode-math} % this also loads fontspec
  \defaultfontfeatures{Scale=MatchLowercase}
  \defaultfontfeatures[\rmfamily]{Ligatures=TeX,Scale=1}
\fi
\usepackage{lmodern}
\ifPDFTeX\else
  % xetex/luatex font selection
\fi
% Use upquote if available, for straight quotes in verbatim environments
\IfFileExists{upquote.sty}{\usepackage{upquote}}{}
\IfFileExists{microtype.sty}{% use microtype if available
  \usepackage[]{microtype}
  \UseMicrotypeSet[protrusion]{basicmath} % disable protrusion for tt fonts
}{}
\makeatletter
\@ifundefined{KOMAClassName}{% if non-KOMA class
  \IfFileExists{parskip.sty}{%
    \usepackage{parskip}
  }{% else
    \setlength{\parindent}{0pt}
    \setlength{\parskip}{6pt plus 2pt minus 1pt}}
}{% if KOMA class
  \KOMAoptions{parskip=half}}
\makeatother
\usepackage{xcolor}
\usepackage[margin=0.75in]{geometry}
\usepackage{longtable,booktabs,array}
\usepackage{calc} % for calculating minipage widths
% Correct order of tables after \paragraph or \subparagraph
\usepackage{etoolbox}
\makeatletter
\patchcmd\longtable{\par}{\if@noskipsec\mbox{}\fi\par}{}{}
\makeatother
% Allow footnotes in longtable head/foot
\IfFileExists{footnotehyper.sty}{\usepackage{footnotehyper}}{\usepackage{footnote}}
\makesavenoteenv{longtable}
\usepackage{graphicx}
\makeatletter
\def\maxwidth{\ifdim\Gin@nat@width>\linewidth\linewidth\else\Gin@nat@width\fi}
\def\maxheight{\ifdim\Gin@nat@height>\textheight\textheight\else\Gin@nat@height\fi}
\makeatother
% Scale images if necessary, so that they will not overflow the page
% margins by default, and it is still possible to overwrite the defaults
% using explicit options in \includegraphics[width, height, ...]{}
\setkeys{Gin}{width=\maxwidth,height=\maxheight,keepaspectratio}
% Set default figure placement to htbp
\makeatletter
\def\fps@figure{htbp}
\makeatother
\usepackage{svg}
\setlength{\emergencystretch}{3em} % prevent overfull lines
\providecommand{\tightlist}{%
  \setlength{\itemsep}{0pt}\setlength{\parskip}{0pt}}
\setcounter{secnumdepth}{-\maxdimen} % remove section numbering
\ifLuaTeX
\usepackage[bidi=basic]{babel}
\else
\usepackage[bidi=default]{babel}
\fi
\babelprovide[main,import]{english}
% get rid of language-specific shorthands (see #6817):
\let\LanguageShortHands\languageshorthands
\def\languageshorthands#1{}
\ifLuaTeX
  \usepackage{selnolig}  % disable illegal ligatures
\fi
\IfFileExists{bookmark.sty}{\usepackage{bookmark}}{\usepackage{hyperref}}
\IfFileExists{xurl.sty}{\usepackage{xurl}}{} % add URL line breaks if available
\urlstyle{same}
\hypersetup{
  pdflang={en},
  hidelinks,
  pdfcreator={LaTeX via pandoc}}

\author{}
\date{}

\begin{document}

\hypertarget{the-theory-of-everything}{%
\section{The Theory of Everything}\label{the-theory-of-everything}}

One Axiom. One Field. All of Physics.

Richard Kent Gates

mail@richardkentgates.com

September 16, 2026

\hypertarget{the-axiom}{%
\subsection{1. The Axiom}\label{the-axiom}}

Distance necessitates time.\\
Time is information.\\
Information cannot be lost.

This single axiom generates quantum mechanics (information = bits =
quantum states), general relativity (distance = time = spacetime),
thermodynamics (information preserved = entropy arrow), and the scalar
field \textbackslash(G(t)\textbackslash) (the mathematical expression of
the axiom).

\hypertarget{the-field}{%
\subsection{2. The Field}\label{the-field}}

The scalar field:

\textbackslash{[}G(t) = 10\^{}\{-3\} +
10\^{}\{-4\}\textbackslash cos(2\textbackslash pi t)\textbackslash{]}

The field IS the quantum vacuum. The quantum vacuum IS spacetime. They
were never separate. They were always the same thing described at
different scales.

\begin{itemize}
\tightlist
\item
  \textbf{GR sees:} spacetime curvature (large-scale field gradient)
\item
  \textbf{QM sees:} vacuum fluctuations (small-scale field oscillation)
\item
  \textbf{Both see:} the same field, the same law, all scales
\end{itemize}

\begin{figure}
\centering
\includesvg{fig-multi-scale-field.svg}
\caption{Figure 1: The same scalar field G(t) at every scale. QM sees
vacuum fluctuations (high-frequency oscillation). GR sees spacetime
curvature (low-frequency gradient). The cosmological constant problem
dissolves.}
\end{figure}

\hypertarget{the-correction-law-the-bridge}{%
\subsection{3. The Correction Law (The
Bridge)}\label{the-correction-law-the-bridge}}

\textbackslash{[}X\_\{\textbackslash text\{eff\}\} = X\_0
\textbackslash left(1 + c\_X \textbackslash,
G(t)\textbackslash right)\textbackslash{]}

This single law applies to ALL sectors:

\begin{longtable}[]{@{}llll@{}}
\toprule\noalign{}
Sector & \textbackslash(X\_\{\textbackslash text\{eff\}\}\textbackslash)
& \textbackslash(c\_X\textbackslash) & Domain \\
\midrule\noalign{}
\endhead
\bottomrule\noalign{}
\endlastfoot
Decay &
\textbackslash(\textbackslash lambda\_\{\textbackslash text\{eff\}\}\textbackslash)
& \textbackslash(Q\_\textbackslash alpha d\_\textbackslash alpha +
Q\_\textbackslash mu d\_\textbackslash mu\textbackslash) & QM
(nuclear) \\
Gravity &
\textbackslash(g\_\{\textbackslash text\{eff\}\}\textbackslash) &
\textbackslash(\textbackslash sum f\_i(d\_\textbackslash alpha +
\textbackslash beta\_i d\_\textbackslash mu)\textbackslash) & GR
(macroscopic) \\
Inertia &
\textbackslash(a\_\{\textbackslash text\{eff\}\}\textbackslash) & same
as time gradient & GR (macroscopic) \\
Fine structure &
\textbackslash(\textbackslash alpha(\textbackslash phi)\textbackslash) &
\textbackslash(d\_\textbackslash alpha\textbackslash) & QM + GR
(fundamental) \\
Mass ratio &
\textbackslash(\textbackslash mu(\textbackslash phi)\textbackslash) &
\textbackslash(d\_\textbackslash mu\textbackslash) & QM + GR
(fundamental) \\
\end{longtable}

The coupling constants are not free parameters. They are derived from QM
sensitivity coefficients, material composition, and two fundamental
dilaton couplings
(\textbackslash(d\_\textbackslash alpha\textbackslash),
\textbackslash(d\_\textbackslash mu\textbackslash)). Two free
parameters. All sectors. All scales.

\begin{figure}
\centering
\includesvg{fig-correction-law-hub.svg}
\caption{Figure 2: The unified correction law as a hub diagram. One
equation, two parameters, five sectors.}
\end{figure}

\hypertarget{what-the-field-does}{%
\subsection{4. What the Field Does}\label{what-the-field-does}}

\begin{itemize}
\tightlist
\item
  Drives expansion (\textbackslash(H\^{}2 \textbackslash sim
  \textbackslash rho\_\{\textbackslash text\{field\}\}\textbackslash))
\item
  Produces radiation (oscillation at large amplitude)
\item
  Produces matter (vibration condensing into relics)
\item
  Produces accelerated expansion (field at small amplitude)
\item
  Modifies decay rates
  (\textbackslash(\textbackslash lambda\_\{\textbackslash text\{eff\}\}
  = \textbackslash lambda\_0(1 + \textbackslash alpha G)\textbackslash))
\item
  Modifies the time gradient
  (\textbackslash(g\_\{\textbackslash text\{eff\}\} = g\_0(1 +
  \textbackslash gamma G)\textbackslash))
\item
  Prevents singularities (\textbackslash(G \textbackslash to
  \textbackslash gamma\^{}\{-1\}\textbackslash) stops collapse)
\item
  Preserves information (no singularities = no information loss)
\end{itemize}

\hypertarget{the-quantum-vacuum-is-spacetime}{%
\subsection{5. The Quantum Vacuum Is
Spacetime}\label{the-quantum-vacuum-is-spacetime}}

The cosmological constant problem (10\textsuperscript{122} orders of
magnitude discrepancy) dissolves: the QM vacuum energy and the observed
accelerated expansion are NOT different things. They are the SAME field
measured at different scales. QM sees the field\textquotesingle s
zero-point oscillation (all modes). GR sees the field\textquotesingle s
macroscopic average (background). The discrepancy is an artifact of
treating them as separate.

\begin{itemize}
\tightlist
\item
  Virtual particles = field fluctuations
\item
  Vacuum energy = field zero-point energy
\item
  Casimir effect = field boundary conditions
\item
  Lamb shift = field coupling to electron
\end{itemize}

\hypertarget{the-measurement-problem-dissolves}{%
\subsection{6. The Measurement Problem
Dissolves}\label{the-measurement-problem-dissolves}}

Superposition = field in undetermined state. Measurement = information
resolves. Collapse = the field\textquotesingle s state becomes
determined by interaction. Not collapse. Information resolution. The
axiom (information cannot be lost) forces the field\textquotesingle s
state to become definite when it interacts. No observer needed. No
consciousness needed. No magic.

\hypertarget{the-entropy-arrow-emerges}{%
\subsection{7. The Entropy Arrow
Emerges}\label{the-entropy-arrow-emerges}}

The field\textquotesingle s fundamental equations are time-reversible
(\textbackslash(dS/dt = 0\textbackslash)). But coarse-graining
(incomplete description) produces \textbackslash(dS/dt \textgreater{}
0\textbackslash). The entropy arrow is not a law. It is a consequence of
information preservation plus incomplete description. The arrow of time
= the direction of information resolution.

\hypertarget{e-mcuxb2-backward}{%
\subsection{8. E = mc² Backward}\label{e-mcuxb2-backward}}

Forward: matter releases energy. Backward: energy condenses into matter.

\begin{itemize}
\tightlist
\item
  Fast vibration (high amplitude) = energy = radiation
\item
  Slow vibration (low amplitude) = matter = relics
\item
  The field damps (Hubble friction): radiation era → matter era →
  expansion era
\end{itemize}

Matter is not a substance. Matter is the field vibration slowed to the
point of condensation. \textbackslash(c\^{}2\textbackslash) is the
field\textquotesingle s maximum vibration speed --- the speed of light =
the speed of information = the speed of state changes.

\hypertarget{status-of-each-problem-under-the-field-model}{%
\subsection{9. Status of Each Problem Under the Field
Model}\label{status-of-each-problem-under-the-field-model}}

The single field and the single correction law are posited in Section 3.
For each problem below, this table states what the companion papers in
this body of work actually demonstrate, and what remains a hypothesis.
Where a row is marked \emph{Hypothesis}, no claim of resolution is made
here; the row records the target the framework is being tested against.

\begin{longtable}[]{@{}lll@{}}
\toprule\noalign{}
Problem & Status under the field model & Basis \\
\midrule\noalign{}
\endhead
\bottomrule\noalign{}
\endlastfoot
GR/QM unification & \emph{Hypothesis.} The correction law is written to
span both domains with two free parameters, but no computation in this
corpus derives quantum behaviour from it. & Posited, Section 3 \\
Cosmological constant & \emph{Hypothesis.} The vacuum energy and
observed dark energy are identified as the same field at different
scales. This is a re-identification; no mechanism producing the ratio is
computed. & Posited, Section 5 \\
Measurement problem & \emph{Hypothesis.} Collapse is recast as
information resolution, following from the axiom that information cannot
be lost. Stated as an argument, not derived. & Section 6 \\
Entropy arrow & Stated. The field equations are time-reversible, giving
{\textbackslash(dS/dt = 0\textbackslash)} fundamentally; the second law
is recovered in the coarse-grained description, so the arrow follows
from information preservation plus incomplete description. &
\texttt{foundation\_proof.py}, \texttt{theory\_of\_everything.py} \\
Black hole singularity & Derived. The time-gradient field
{\textbackslash(G = d(d\textbackslash tau/dt)/d(\textbackslash ln
r)\textbackslash)} diverges as the Schwarzschild radius is approached.
Since {\textbackslash(G\_\{\textbackslash mathrm\{eff\}\} = G\_\{0\}(1 -
c\_\{g\}G)\textbackslash)} is regulated by that field, it reaches zero
first: freeze-out at {\textbackslash(r/r\_\{s\} = (1 +
\textbackslash sqrt\{1+c\_\{g\}\^{}\{2\}\})/2\textbackslash)}, outside
the horizon for every {\textbackslash(c\_\{g\} \textgreater{}
0\textbackslash)}. No horizon forms. & \texttt{blackhole\_lifecycle.py},
9/9 checks \\
Gravity as a time gradient & Measured. The law {\textbackslash(a =
c\^{}\{2\}\textbackslash,d(d\textbackslash tau/dt)/dr\textbackslash)}
predicts the rate difference between two clocks with no fitted
parameter. At 450 m the GR-violation parameter is consistent with the
zero-correction limit at 0.15σ; at 457 km chronometric and geodetic
determinations of the same time-density difference agree to 0.85σ; at
1415 km the residual after the potential is removed is consistent with
zero. & \texttt{chronometric\_levelling.py}; Takamoto et al. 2020,
arXiv:2309.14953, Grose et al. 2015 \\
Rotation curve structure & Fitted. Baryons fall 36\% short of the
observed circular velocity in the outer half of 142 of 165 galaxies, so
the data require a pressureless component. The gradient law with one
free parameter per galaxy reduces the median residual from 33.5 to 10.1
km/s, and the pressureless component is preferred in 165 of 165. &
\texttt{rotation\_curves\_sparc.py}; Lelli et al. 2019 (SPARC, CC BY
4.0) \\
Big Bang singularity & Stated as a prohibition. A singularity requires
infinite information in zero volume, violating the Bekenstein bound; the
quartic potential bounds the field amplitude in both time directions.
Argued, not computed. & \texttt{foundation\_proof.py}, 8/8 checks \\
Hierarchy problem & Partial. The field is dimensionless and couples
conformally, so loop corrections are logarithmic rather than quadratic.
Dimensional argument, not a renormalization calculation. &
\texttt{renormalization\_proof.py}, 7/7 checks \\
Hubble tension & \emph{Hypothesis.} Environment-dependent field
amplitude is offered as an account. No prediction for the
tension\textquotesingle s magnitude is derived in this corpus. &
Posited \\
Correlated decay anomalies & Documented. Five independent datasets show
fractional deviations clustering in the
\textbackslash(10\^{}\{-4\}\textbackslash) to
\textbackslash(10\^{}\{-3\}\textbackslash) band under a single equation,
with fitted couplings within an order of magnitude across isotopes. &
\texttt{time-gradient-field-model.html},
\texttt{time\_gradient\_verification.html} \\
Statistical recoverability & Numerically demonstrated. The field is
reconstructible from noisy multi-sector measurement; empirical variance
is 2.73 times the Cramér--Rao bound, with all three estimators
converging to the same error floor. &
\texttt{signal\_processing\_proof.py},
\texttt{statistical-verification.html} \\
\end{longtable}

Two of the ten rows rest on numerical demonstration in this corpus. The
remainder are either stated arguments or explicit hypotheses, marked
above. The clustering of decay anomalies and the statistical
recoverability of the field are the results this body of work
establishes; the unification and cosmological claims are the targets it
is measured against.

\hypertarget{conclusion}{%
\subsection{10. Conclusion}\label{conclusion}}

One axiom: distance → time → information → preserved.

One field: \textbackslash(G(t)\textbackslash).

One law: \textbackslash(X\_\{\textbackslash text\{eff\}\} = X\_0(1 +
c\_X G)\textbackslash).

All of physics.

The Theory of Everything is not a new force or a new particle. It is a
new understanding of what already exists: the field IS. The field
changes. Change IS time. Time IS information. Information cannot be
lost. The field cannot end. There is only the field, changing, forever.

\hypertarget{ai-research-collaboration-disclosure}{%
\subsection{AI Research Collaboration
Disclosure}\label{ai-research-collaboration-disclosure}}

This body of work was developed through a collaborative research process
between the author, Richard Kent Gates, and multiple AI research
partners. The author provides full transparency on this process.

\textbf{Author\textquotesingle s Role (Richard Kent Gates):}

\begin{itemize}
\tightlist
\item
  Conceived all core theories, conceptual frameworks, hypotheses, and
  research direction
\item
  Defined the Time-Gradient Field Model, the unified scalar field G(t),
  the distance → time → information → G(t) framework, and all
  theoretical architecture
\item
  Provided physical intuition, biological reasoning, and domain
  expertise that guided every theoretical decision
\item
  Reviewed, validated, and approved all mathematical formalisms before
  inclusion
\item
  Made all final judgments on theoretical coherence and physical
  interpretation
\end{itemize}

\textbf{AI Research Partners:}

\begin{itemize}
\tightlist
\item
  \textbf{Mimo V2.5 (opencode platform)} --- Primary mathematical
  formalization partner. Assisted in translating conceptual physics into
  mathematical notation, generating numerical proof scripts, compiling
  LaTeX documents, and building the website infrastructure.
\item
  \textbf{Microsoft Copilot (browser-based)} --- Assisted in literature
  review, dataset gathering, cross-referencing empirical results (DESI,
  BNL, PTB, MICROSCOPE), and verifying citation accuracy.
\item
  \textbf{Google Gemini (browser-based)} --- Assisted in conceptual
  exploration, thought experimentation, and early-stage hypothesis
  refinement.
\item
  \textbf{Space Bunny (OpenCode)} --- Verification and provenance
  auditing. Read every paper and proof script in the repository and
  traced each reported figure back to the script and line that produces
  it. Independently recomputed the Fisher information matrices, the
  Cramér--Rao bounds, and the parameter errors from the models as
  stated. Resolved the origin of the Fisher information values quoted in
  the unified and Theory-of-Everything papers, reconstructed the
  two-parameter Fisher matrix in the statistical verification paper from
  its sampling grid, and identified a mislabeled citation and a
  statistical criterion stated more narrowly than the reported results
  supported. No theoretical claim, numerical result, or interpretive
  judgment in this body of work originated with this tool.
\end{itemize}

\textbf{Nature of AI Involvement:}

The AI tools functioned as research assistants --- analogous to graduate
students or technical collaborators who help formalize, compute, and
organize ideas that originate from the principal investigator. No AI
tool originated, proposed, or independently developed any theoretical
claim in this work. All physical insights, theoretical innovations, and
interpretive judgments are the author\textquotesingle s own.

\hypertarget{references}{%
\subsection{References}\label{references}}

{[}1{]} Weinberg, S. "The Cosmological Constant Problem." \emph{Reviews
of Modern Physics}, 61(1), 1--23, 1989.

{[}2{]} Casimir, H.B.G. "On the Attraction Between Two Perfectly
Conducting Plates." \emph{Proceedings of the Koninklijke Nederlandse
Akademie van Wetenschappen}, 51, 793--795, 1948.

{[}3{]} Lamb, W.E. \& Retherford, R.C. "Fine Structure of the Hydrogen
Atom by a Microwave Method." \emph{Physical Review}, 72(3), 241--243,
1947.

{[}4{]} Einstein, A. "Ist die Trägheit eines Körpers von seinem
Energieinhalt abhängig?" \emph{Annalen der Physik}, 323(13), 639--641,
1905.

{[}5{]} Hawking, S.W. "Breakdown of Predictability in Gravitational
Collapse." \emph{Physical Review D}, 14(10), 2460--2473, 1976.

{[}6{]} Penrose, R. "Gravitational Collapse and Space-Time
Singularities." \emph{Physical Review Letters}, 14(3), 57--59, 1965.

{[}7{]} Riess, A.G. et al. "A Comprehensive Measurement of the Local
Value of the Hubble Constant." \emph{The Astrophysical Journal Letters},
934(1), L7, 2022.

{[}8{]} Giudice, G.F. "Naturalizing Hierarchy." \emph{arXiv:0801.2562},
2008.

{[}9{]} Hawking, S.W. "Particle Creation by Black Holes."
\emph{Communications in Mathematical Physics}, 43(3), 199--220, 1975.

{[}10{]} Wheeler, J.A. \& Zurek, W.H. \emph{Quantum Theory and
Measurement}. Princeton University Press, 1983.

\href{https://richardkentgates.com/}{← richardkentgates.com}

\end{document}
